\documentclass[../main.tex]{subfiles}
\ifSubfilesClassLoaded{\setcounter{chapter}{7}}{}
\begin{document}

\chapter{Berkas Teks, Format Output Lanjut, dan Penanganan Kesalahan}

\begin{subcpmk}
  \item Sub-CPMK 5.1: Membaca dan menulis berkas teks dengan konteks manajer (\texttt{with}) dan penanganan kesalahan dasar
  \item Sub-CPMK 5.2: Memuat dan menyimpan data terstruktur JSON serta menerapkan \texttt{try}/\texttt{except}/\texttt{else}/\texttt{finally} secara tepat
\end{subcpmk}

\SesiRPS{9}{5.1 dan 5.2}{$\sim$170 menit}{CPMK-5}

\section{Sarana dan prasarana}

Python 3.10+, editor, terminal; berkas contoh teks/JSON di folder kerja.

\section{Prosedur kerja}

\begin{enumerate}[leftmargin=*]
  \item Latih format \texttt{f-string} dan penulisan berkas teks.
  \item Baca/tulis JSON kecil; tangani pengecualian umum:\newline
  \texttt{FileNotFoundError}, \texttt{JSONDecodeError}.
  \item Dokumentasikan jalur berkas relatif terhadap skrip.
\end{enumerate}

\section{Format Output Lanjut}

Sebelum membahas berkas, mari pelajari cara memformat output dengan baik.

\subsection{f-string Lanjut dan \texttt{str.format()}}

Python menyediakan beberapa cara memformat string (Tutorial Python §7.1):

\begin{lstlisting}[caption=Format string: f-string dan format() (Tutorial Python ss7.1)]
# f-string dengan format specification
nama = "Aulia"
skor = 87.456
print(f"Nama  : {nama:<15}")     # rata kiri lebar 15
print(f"Skor  : {skor:>8.2f}")   # rata kanan, 2 desimal
print(f"Persen: {skor/100:.1%}") # format persen

# str.format() -- alternatif f-string
print("Nama: {:15} | Skor: {:6.2f}".format(nama, skor))

# repr() vs str() untuk debugging
print(f"{nama!r}")   # 'Aulia' -- dengan tanda kutip
print(f"{nama!s}")   # Aulia  -- tanpa tanda kutip
\end{lstlisting}

\section{Membaca dan Menulis Berkas}

Gunakan \texttt{open} dengan mode \texttt{r}, \texttt{w}, \texttt{a}. Disarankan memakai konteks \keyword{with} agar berkas tertutup otomatis (Tutorial Python §7.2):

\begin{lstlisting}[caption=Menulis dan membaca berkas]
# Menulis ke berkas
with open("data.txt", "w", encoding="utf-8") as f:
    f.write("Baris pertama\n")
    f.write("Baris kedua\n")

# Membaca seluruh isi
with open("data.txt", "r", encoding="utf-8") as f:
    isi = f.read()
print(isi)
\end{lstlisting}

\section{Metode Objek Berkas Lanjut}

Objek berkas memiliki beberapa metode yang berguna (Tutorial Python §7.2.1):

\begin{lstlisting}[caption=Metode berkas: readline readlines seek tell]
with open("data.txt", "w", encoding="utf-8") as f:
    for i in range(1, 6):
        f.write(f"Baris {i}: data penting\n")

with open("data.txt", "r", encoding="utf-8") as f:
    baris1 = f.readline()          # baca satu baris
    print(repr(baris1))            # 'Baris 1: data penting\n'

    semua = f.readlines()          # baca semua sisa sebagai list
    print(len(semua))              # 4 (baris 2--5)

# seek dan tell: navigasi posisi berkas
with open("data.txt", "rb") as f:    # buka mode binary untuk seek
    print(f.tell())      # 0 -- posisi awal
    f.seek(10)           # pindah ke byte ke-10
    print(f.tell())      # 10
    f.seek(0)            # kembali ke awal
    print(f.read(5))     # baca 5 byte pertama
\end{lstlisting}

\section{Iterasi Baris}

Cara paling efisien membaca berkas baris per baris (hemat memori untuk file besar):

\begin{lstlisting}[caption=Iterasi baris -- cara paling efisien]
with open("data.txt", "r", encoding="utf-8") as f:
    for nomor, baris in enumerate(f, start=1):
        print(f"{nomor:3}: {baris}", end="")
\end{lstlisting}

\section{Menyimpan Data Terstruktur dengan JSON}

Modul \texttt{json} memudahkan penyimpanan dan pembacaan data terstruktur\newline (Tutorial Python §7.2.2):

\begin{lstlisting}[caption=Membaca dan menulis JSON]
import json

# Data Python -> JSON string
data = {
    "mahasiswa": [
        {"nim": "12345", "nama": "Andi", "ipk": 3.5},
        {"nim": "12346", "nama": "Budi", "ipk": 3.8},
    ],
    "semester": 1
}

# Simpan ke berkas JSON
with open("data.json", "w", encoding="utf-8") as f:
    json.dump(data, f, indent=2, ensure_ascii=False)

# Baca dari berkas JSON
with open("data.json", "r", encoding="utf-8") as f:
    data_terbaca = json.load(f)

for mhs in data_terbaca["mahasiswa"]:
    print(f"{mhs['nim']}: {mhs['nama']} (IPK {mhs['ipk']})")

# Konversi ke/dari string JSON
teks_json = json.dumps(data, indent=2)
obj_kembali = json.loads(teks_json)
\end{lstlisting}

\section{Penanganan Kesalahan Lanjut: \texttt{try} / \texttt{except} / \texttt{else} / \texttt{finally}}

Python menggunakan \textit{exception} untuk menangani kesalahan runtime (Tutorial Python §8):

\begin{lstlisting}[caption=Pola try/except/else/finally lengkap]
def baca_angka(prompt: str) -> int:
    """Membaca angka bulat dari pengguna dengan validasi."""
    while True:
        try:
            nilai = int(input(prompt))
        except ValueError:
            print("Masukan tidak valid. Masukkan angka bulat.")
        else:
            # Dieksekusi jika try BERHASIL (tidak ada exception)
            return nilai

# Beberapa jenis exception sekaligus
def buka_berkas(nama: str) -> str:
    try:
        with open(nama, "r", encoding="utf-8") as f:
            return f.read()
    except FileNotFoundError:
        print(f"Berkas '{nama}' tidak ditemukan.")
        return ""
    except PermissionError:
        print(f"Tidak ada izin membaca '{nama}'.")
        return ""
    except OSError as e:
        print(f"Kesalahan OS: {e}")
        return ""
    finally:
        # SELALU dieksekusi, terlepas dari hasil
        print(f"Percobaan membuka '{nama}' selesai.")
\end{lstlisting}

\begin{lstlisting}[caption=Raise exception dan exception kustom]
def bagi(a: float, b: float) -> float:
    if b == 0:
        raise ValueError("Pembagi tidak boleh nol!")
    return a / b

try:
    hasil = bagi(10, 0)
except ValueError as e:
    print(f"Error: {e}")   # Error: Pembagi tidak boleh nol!
\end{lstlisting}

\begin{konsep}
\textbf{Alur \texttt{try}/\texttt{except}/\texttt{else}/\texttt{finally}:}
\begin{itemize}[leftmargin=*]
  \item \textbf{\texttt{try}:} Kode yang mungkin memunculkan exception
  \item \textbf{\texttt{except}:} Menangani exception jika terjadi
  \item \textbf{\texttt{else}:} Dieksekusi \textit{hanya jika} try berhasil tanpa exception
  \item \textbf{\texttt{finally}:} \textit{Selalu} dieksekusi (untuk cleanup, menutup koneksi, dll.)
\end{itemize}
\end{konsep}

\section{Studi Kasus dan Contoh Kode}

Berikut adalah contoh-contoh program Python yang mengimplementasikan konsep-konsep pada modul ini.

\subsection{Kode: \texttt{baca\_tulis.py}}
\begin{lstlisting}[language=Python,caption={\texttt{baca\_tulis.py}}]
"""Menulis dan membaca file teks dengan penanganan kesalahan dasar."""


def main() -> None:
    nama_file = "data_mahasiswa.txt"
    
    # Menulis ke file
    try:
        with open(nama_file, "w", encoding="utf-8") as f:
            f.write("Budi, 3.8\n")
            f.write("Siti, 3.9\n")
        print(f"Data berhasil ditulis ke {nama_file}.")
    except IOError as e:
        print(f"Gagal menulis file: {e}")

    # Membaca dari file
    try:
        with open(nama_file, "r", encoding="utf-8") as f:
            print("\nIsi file:")
            for baris in f:
                print(baris.strip())
    except FileNotFoundError:
        print(f"File {nama_file} tidak ditemukan.")

if __name__ == "__main__":
    main()
\end{lstlisting}

\subsection{Kode: \texttt{olah\_json.py}}
\begin{lstlisting}[language=Python,caption={\texttt{olah\_json.py}}]
"""Menyimpan dan memuat data JSON."""
import json


def main() -> None:
    data = {
        "mata_kuliah": "Praktikum Pemrograman",
        "peserta": ["Andi", "Budi", "Citra"]
    }
    file_json = "data.json"
    
    with open(file_json, "w", encoding="utf-8") as f:
        json.dump(data, f, indent=4)
        print("Data JSON telah disimpan.")
        
    with open(file_json, "r", encoding="utf-8") as f:
        data_baca = json.load(f)
        print("Data JSON dibaca:", data_baca["mata_kuliah"])

if __name__ == "__main__":
    main()
\end{lstlisting}

\begin{aktivitas}
  \item Buat program yang menyalin isi satu berkas teks ke berkas lain (tanpa \texttt{shutil}, gunakan loop baca/tulis).
  \item Hitung jumlah baris, kata, dan karakter pada berkas teks dummy.\newline Gunakan \texttt{enumerate}.
  \item Simpan data mahasiswa (list of dict) ke berkas JSON, kemudian baca kembali dan tampilkan rekap.
  \item Tambahkan penanganan \texttt{try}/\texttt{except}/\texttt{finally} ke program pembacaan berkas; pastikan pesan kesalahan informatif.
  \item Simpan log aktivitas (timestamp + pesan) dengan mode \texttt{a} dan baca log tersebut dengan format yang rapi.
\end{aktivitas}

\begin{latihan}
  \item Mengapa \texttt{with} disarankan dibanding \texttt{open}/\texttt{close} manual?
  \item Apa beda mode \texttt{w} dan \texttt{a}? Apa yang terjadi jika berkas belum ada?
  \item Kapan menggunakan \texttt{readline()} vs \texttt{readlines()} vs iterasi langsung?
  \item Apa perbedaan \texttt{json.dumps()} dan \texttt{json.dump()}? Kapan masing-masing digunakan?
  \item Apa yang dilakukan blok \texttt{else} pada \texttt{try}? Mengapa ini berguna?
  \item Refleksi: kesalahan apa yang sering muncul saat path relatif di Windows vs Linux?
\end{latihan}

\begin{asesmen}
\textbf{Cek pemahaman}

\begin{enumerate}
  \item Mode \texttt{open} untuk menambah di akhir berkas tanpa menghapus isi lama:
  \begin{enumerate}
    \item \texttt{"r"}
    \item \texttt{"w"}
    \item \texttt{"a"}
    \item \texttt{"x"}
  \end{enumerate}
  \item Fungsi \texttt{json.load(f)} digunakan untuk:
  \begin{enumerate}
    \item Mengubah objek Python menjadi JSON string
    \item Membaca JSON dari objek berkas dan mengubahnya ke Python
    \item Menampilkan isi JSON dengan format
    \item Menyimpan JSON ke berkas
  \end{enumerate}
  \item Blok \texttt{finally} dalam \texttt{try/except} dieksekusi:
  \begin{enumerate}
    \item Hanya saat tidak ada exception
    \item Hanya saat ada exception
    \item Selalu, terlepas dari ada tidaknya exception
    \item Hanya jika ada blok \texttt{else}
  \end{enumerate}
\end{enumerate}

\textbf{Tugas}: program \texttt{rekap\_nilai.py} yang membaca data mahasiswa dari berkas JSON, menghitung statistik (rata-rata, nilai tertinggi, nilai terendah), dan menyimpan hasil ke berkas teks dengan format tabel. Gunakan \texttt{try/except} untuk menangani berkas yang tidak ditemukan.
\end{asesmen}

\begin{checklist}
  \item Saya dapat membaca seluruh berkas atau per baris (\texttt{read}, \texttt{readline}, \texttt{readlines})
  \item Saya dapat menulis dan menambah berkas dengan mode yang tepat
  \item Saya memahami penggunaan \texttt{seek()} dan \texttt{tell()} untuk navigasi berkas
  \item Saya dapat menyimpan dan membaca data JSON dengan \texttt{json.dump}/\texttt{json.load}
  \item Saya dapat menggunakan \texttt{try}/\texttt{except}/\texttt{else}/\texttt{finally}
  \item Saya dapat menggunakan \texttt{raise} untuk memunculkan exception
  \item Saya memahami pentingnya \texttt{encoding="utf-8"}
  \item Saya dapat memformat output dengan f-string dan format specification
\end{checklist}

\section{Referensi sesi}

\begin{itemize}[leftmargin=*]
  \item \textbf{[19]} Real Python --- berkas, I/O, dan topik terkait.
  \item \textbf{[8, 9]} \textit{The Python Tutorial}, masukan/keluaran dan kesalahan.
\end{itemize}

\begin{rangkuman}
Akses berkas memerlukan manajemen sumber daya yang benar --- \texttt{with} membantu menutup berkas meskipun terjadi error. JSON adalah format pertukaran data yang ringkas untuk menyimpan struktur data Python. Penanganan exception yang tepat membuat program lebih robust dan ramah pengguna.

\textbf{Referensi:} {\small Python Tutorial §7, §8 \url{https://docs.python.org/3/tutorial/inputoutput.html} dan \url{https://docs.python.org/3/tutorial/errors.html}}

\textbf{Kata kunci:} \code{open}, \code{with}, \code{read}, \code{readline}, \code{seek}, \code{json}, \code{try}, \code{except}, \code{else}, \code{finally}, \code{raise}
\end{rangkuman}

\end{document}
